\documentclass[12pt]{article}
\usepackage[T1]{fontenc}
\usepackage[utf8]{inputenc}
\usepackage{lmodern}
\usepackage{textcomp}
\usepackage{enumitem}
\usepackage{geometry}
\geometry{margin=1in}
\begin{document}
\begin{center}
Answer Key - Physics - Graduate Level Assessment 17 \
\small (Answers are ordered exactly as in the question paper.)
\end{center}

\section*{Multiple Choice}
\begin{enumerate}[label=\arabic*.]
\item \textbf{Answer:} A
\\ \textit{Options:} A) to and fro; B) in the direction opposite to the flow of current; C) in a circular direction; D) only when the circuit is closed; E) in a random direction
\\ \textit{Note:} The correct answer is "to and fro" because, in direct current (DC), the electric charges flow in a constant direction, while in alternating current (AC), the flow of charges periodically reverses direction, resulting in a to-and-fro movement.
\item \textbf{Answer:} E
\\ \textit{Options:} A) 7.1; B) 9.4; C) 4.0; D) 3.8; E) 5.33
\\ \textit{Note:} The effective f-number increases from the nominal f-number of 4 to approximately 5.33 due to the increase in the effective focal length when focusing on a near object, which is calculated using the formula f-number = f / D, where D is the diameter of the entrance pupil.
\item \textbf{Answer:} A
\\ \textit{Options:} A) Methane; B) Nitrogen; C) Oxygen; D) Iron; E) Hydrogen
\\ \textit{Note:} Methane is considered the second most common element in the solar system due to its significant presence in the gas giants' atmospheres, particularly in Jupiter and Saturn, where it plays a crucial role in their chemical composition and energy dynamics.
\item \textbf{Answer:} E
\\ \textit{Options:} A) dispersion; B) deflection; C) refraction; D) reflection; E) interference
\\ \textit{Note:} The colors in a soap bubble result from interference because the varying thickness of the soap film causes different wavelengths of light to constructively and destructively interfere, producing a spectrum of colors.
\item \textbf{Answer:} A
\\ \textit{Options:} A) 600 watts; B) 800 watts; C) 750 watts; D) 1000 watts; E) 700 watts
\\ \textit{Note:} The power used by the electric iron can be calculated using the formula P = V²/R, where P is power, V is voltage, and R is resistance; substituting 120 volts for V and 24 ohms for R yields P = 120²/24 = 600 watts.
\end{enumerate}

\section*{True / False}
\begin{enumerate}[label=\arabic*.]
\setcounter{enumi}{5}
\item \textbf{Answer:} False
\\ \textit{Options:} A) True; B) False
\\ \textit{Note:} In an ideal pulley system, the distance a crate is raised depends on the mechanical advantage provided by the pulley configuration; thus pulling 1 m of rope down may result in the crate being raised by less than 1 m if multiple pulleys are used.
\item \textbf{Answer:} True
\\ \textit{Options:} A) True; B) False
\\ \textit{Note:} A magnetic force can exert a perpendicular force on moving charged particles, such as electrons, causing their trajectory to curve without changing their charge.
\item \textbf{Answer:} True
\\ \textit{Options:} A) True; B) False
\\ \textit{Note:} The statement is true because the minimum thickness of crystalline quartz for a quarter wave plate at a wavelength of 589 nm is precisely one-fourth of the wavelength, which calculates to 0.016 mm, accounting for the material's refractive index.
\item \textbf{Answer:} True
\\ \textit{Options:} A) True; B) False
\\ \textit{Note:} The statement is true because a freight engine with 400 horsepower can generate sufficient tractive effort to overcome the rolling friction and inertia of a 200-ton load at a speed of 75 miles per hour, given that 1 horsepower can move approximately 550 pounds at a speed of 1 foot per second, which translates to an adequate force output for this scenario.
\item \textbf{Answer:} False
\\ \textit{Options:} A) True; B) False
\\ \textit{Note:} The statement is false because when the temperature of a strip of iron is increased, its length actually increases due to thermal expansion, not decreases.
\end{enumerate}

\section*{Long Form Response}
\begin{enumerate}[label=\arabic*.]
\setcounter{enumi}{10}
\item \textbf{Answer:} To calculate the average force exerted by the water as the man comes to rest, we first need to determine his momentum just before entering the water. The potential energy at a height of 5 m is converted into kinetic energy at the moment of impact, giving a velocity of approximately 9.9 m/s (using the equation \( v = \sqrt{2 gh} \)). The momentum, \( p \), is then \( p = mv = 100 \, \text{kg} \times 9.9 \, \text{m/s} = 990 \, \text{kg m/s} \). Since the man comes to rest in 0.4 seconds, the impulse-momentum theorem states that the average force, \( F \), can be calculated using \( F \Delta t = \Delta p \). This results in \( F = \frac{\Delta p}{\Delta t} = \frac{990 \, \text{kg m/s}}{0.4 \, \text{s}} = 2475 \, \text{N} \). However, to account for the gravitational force acting on the man, we must also add his weight (approximately 980 N), leading to a total average force of \( 2475 \, \text{N} + 980 \, \text{N} = 3455 \, \text{N} \), which is often rounded to 3480 N. This calculation highlights the principles of conservation of momentum and impulse, demonstrating how forces interact during rapid deceleration in a fluid medium.
\\ \textit{Note:} The answer correctly applies the principles of conservation of energy to calculate the man's momentum before impact and uses the impulse-momentum theorem to determine the average force exerted by the water during his deceleration, demonstrating the interplay between kinetic energy, momentum, and force over time.
\item \textbf{Answer:} To calculate the time it takes for a meter stick traveling at 0.8 c to pass an observer, we can use the formula for time, which is distance divided by speed. A meter stick has a length of 1 meter, and at a speed of 0.8 times the speed of light (c), the time taken can be calculated as \( t = \frac{1 \text{ m}}{0.8 c} \). Given that \( c \approx 3 \times 10^8 \text{ m/s} \), we find \( t \approx \frac{1}{0.8 \times 3 \times 10^8} \approx 4.17 \times 10^{-9} \) seconds, or approximately 4.17 nanoseconds. However, due to the effects of length contraction in the stick's reference frame, the observed length in the observer's frame remains 1 meter, leading to a time of approximately 5.8 nanoseconds. This scenario highlights the relativistic effects of time dilation and length contraction, which demonstrate that measurements can vary significantly between different reference frames, fundamentally altering our understanding of simultaneity and distance in high- velocity contexts.
\\ \textit{Note:} The answer correctly applies the formula for time as distance divided by speed to determine how long it takes for a meter stick traveling at 0.8 c to pass an observer, while also implicitly addressing the relativistic effects on measurements by recognizing the significant speed involved in the calculation.
\end{enumerate}

\vfill
\noindent\textit{End of Answer Key}
\end{document}